\subsection{环同态}

定义 2. 设 A 和 B 是两个环。从 A 到 B 的一个态射，或同态，是满足以下关系的从 A 到 B 的任意映射 f：

\[
f (x + y) = f (x) + f (y), \quad f (x y) = f (x). f (y), \quad f (1) = 1,\tag{17}
\]

对所有 \(x, y\) （在 A 中）成立。

两个环同态的复合是一个环同态。设 A 和 B 是两个环，f 是 A 到 B 的一个映射；要使 f 成为同构，必要且充分的是它为一个双射同态；在这种情况下， \(\overline{f}^{1}\) 是 B 到 A 的一个同态。环 A 到自身的同态称为 A 的自同态。

设 \(f: A \to B\) 是一个环同态。映射 \(f\) 是 \(A\) 的加法群到 \(B\) 的加法群的一个同态；特别地， \(f(0) = 0\) 和 \(f(-x) = -f(x)\) 对所有 \(x \in A\) 成立。\(f\) 把 \(A\) 的可逆元素映为 \(B\) 的可逆元素，并且 \(f\) 诱导出 \(A\) 的乘法群到 \(B\) 的乘法群中的一个同态。

例子。(1) 设 A 是一个环。立即可以看出，映射 \(n \mapsto n\) . 1 （从 Z 到 A 中）是 Z 到 A 的唯一同态。特别地，Z 的恒等映射是环 Z 的唯一自同态。

特别地，取 A 为加法群 Z 的自同态环（第 3 号，例 VI）。映射 \(n \mapsto n\) . 1 （从 Z 到 A 中）正是由于 Z 中乘法的构造（§ 2，第 6 号）而是 Z 到 A 上的一个同构。

\begin{enumerate}
\def\labelenumi{(\arabic{enumi})}
\setcounter{enumi}{1}
\tightlist
\item
  设 \(a\) 是环 \(A\) 的一个可逆元素。映射 \(x \mapsto axa^{-1}\) 是 \(A\) 的一个自同态，因为
\end{enumerate}

\[
\begin{array}{r l} a (x + y) a ^ {- 1} & = a x a ^ {- 1} + a y a ^ {- 1}, \\ a (x y) a ^ {- 1} & = (a x a ^ {- 1}) (a y a ^ {- 1}). \end{array}
\]

它是双射的，因为关系 \(x' = axa^{-1}\) 等价于 \(x = a^{-1}x'a\) 。因此它是环 A 的一个自同构，称为与 a 相伴的内自同构。

